\chapter{Những Điều Về Nỗ Lực}
\markboth{Những Điều Về Nỗ Lực}{Things a Teacher Wants to Say About Effort}

\section{Cố Gắng Mỗi Ngày}
\begin{truyen}{Cố Gắng Mỗi Ngày}{Try Every Day}
\chuhoa{N}{ỗ lực mạnh nhất không phải lúc nào cũng ồn ào. Nhiều khi nó chỉ là việc hôm nay em vẫn ngồi xuống học thêm một chút.}
\textit{(The strongest effort is not always loud. Often it is simply sitting down and studying a little again today.)}

Cố gắng mỗi ngày giữ cho đường học không đứt.
\textit{(Daily effort keeps the learning road from breaking.)}

Không cần ngày nào cũng làm lớn. Chỉ cần ngày nào cũng còn liên hệ với việc mình phải học.
\textit{(Not every day needs dramatic effort. You only need to keep some living contact with what must be learned.)}

Người cố gắng đều thường đi xa hơn người chỉ bùng lên từng đợt.
\textit{(Those who work steadily often go farther than those who flare up in bursts.)}
\end{truyen}

\begin{baihoc}
Nỗ lực hằng ngày nhỏ nhưng thật thường đáng tin hơn những lần hứa rất lớn rồi tắt nhanh.
\textit{(Small but real daily effort is often more reliable than huge promises that fade quickly.)}
\end{baihoc}

\begin{ghinhoanh}
\textbf{Short English to remember:}

Daily effort keeps learning alive.

\end{ghinhoanh}

\begin{tuhocnhanh}
1. Vì sao nỗ lực mỗi ngày mạnh hơn bùng lên rồi tắt? \textit{Why is daily effort stronger than short bursts?}

2. Hôm nay em có thể giữ liên hệ với việc học bằng việc gì? \textit{How can you stay connected to learning today?}

3. Em thường đứt nhịp ở đâu? \textit{Where do you usually lose your rhythm?}
\end{tuhocnhanh}
\ngancach

\section{Nỗ Lực Nhỏ Tạo Nên Kết Quả Lớn}
\begin{truyen}{Nỗ Lực Nhỏ Tạo Nên Kết Quả Lớn}{Small Efforts Create Large Results}
\chuhoa{K}{ết quả lớn hiếm khi rơi xuống từ một ngày duy nhất. Nó thường được ghép từ rất nhiều bước nhỏ.}
\textit{(Large results rarely fall from one single day. They are usually built from many small steps.)}

Một chút đọc thêm, một chút sửa lỗi, một chút luyện lại.
\textit{(A little more reading, a little more correction, a little more practice.)}

Mỗi phần nhìn riêng thì nhỏ, nhưng cộng lại lâu ngày thì khác biệt.
\textit{(Each part seems small alone, but over time they make a difference.)}

Đừng coi thường nỗ lực nhỏ chỉ vì nó chưa cho em cảm giác chiến thắng ngay.
\textit{(Do not underestimate small effort just because it does not give immediate victory.)}
\end{truyen}

\begin{baihoc}
Cái lớn thường được xây từ việc không bỏ qua cái nhỏ.
\textit{(Big change is often built by not neglecting small work.)}
\end{baihoc}

\begin{ghinhoanh}
\textbf{Short English to remember:}

Small efforts grow into big results over time.

\end{ghinhoanh}

\begin{tuhocnhanh}
1. Vì sao bước nhỏ dễ bị coi thường? \textit{Why are small efforts easily ignored?}

2. Kết quả lớn thường được tạo ra thế nào? \textit{How are large results usually built?}

3. Bước nhỏ nào em đang bỏ qua? \textit{What small effort are you neglecting?}
\end{tuhocnhanh}
\ngancach

\section{Chăm Chỉ Quan Trọng Hơn Tài Năng}
\begin{truyen}{Chăm Chỉ Quan Trọng Hơn Tài Năng}{Hard Work Often Matters More Than Talent}
\chuhoa{T}{ài năng giúp khởi đầu nhẹ hơn. Chăm chỉ giúp em không bị bỏ lại ở đường dài.}
\textit{(Talent may make the start easier. Hard work keeps you from being left behind over the long road.)}

Có những người sáng dạ nhưng lơi tay.
\textit{(Some people are gifted but careless.)}

Cũng có những người không nổi bật ngay, nhưng nhờ chăm mà ngày một chắc hơn.
\textit{(Others are not immediately outstanding, but grow steadily through hard work.)}

Đường học rất công bằng ở chỗ nó thưởng cho người chịu làm việc thật lâu hơn nhiều người tưởng.
\textit{(Learning is fair in this way: it rewards those willing to work longer than many imagine.)}
\end{truyen}

\begin{baihoc}
Tài năng là lợi thế. Chăm chỉ mới là thứ giữ cho lợi thế ấy không bị phí đi.
\textit{(Talent is an advantage. Hard work keeps that advantage from being wasted.)}
\end{baihoc}

\begin{ghinhoanh}
\textbf{Short English to remember:}

Talent starts; hard work sustains.

\end{ghinhoanh}

\begin{tuhocnhanh}
1. Vì sao chăm chỉ có sức mạnh riêng? \textit{Why does hard work have its own power?}

2. Tài năng bị phí đi trong những trường hợp nào? \textit{When does talent get wasted?}

3. Em đang dựa quá nhiều vào điều gì thay vì làm việc thật? \textit{What are you relying on instead of doing the real work?}
\end{tuhocnhanh}
\ngancach

\section{Không Có Con Đường Tắt}
\begin{truyen}{Không Có Con Đường Tắt}{There Is No Real Shortcut}
\chuhoa{A}{i cũng thích con đường nhanh, nhưng học thật thường không đi bằng lối tắt.}
\textit{(Everyone likes the fast road, but real learning usually does not travel by shortcuts.)}

Có thể em học mẹo để qua một bài.
\textit{(You may use tricks to get through one task.)}

Nhưng nếu muốn có nền vững, em vẫn phải đọc, hiểu, luyện và sửa.
\textit{(But if you want a strong foundation, you still need to read, understand, practice, and correct.)}

Con đường tắt thường lấy đi của em chiều sâu mà em chưa kịp thấy mình đang mất.
\textit{(Shortcuts often steal depth before you even notice the loss.)}
\end{truyen}

\begin{baihoc}
Muốn có năng lực thật, sớm muộn em cũng phải đi qua phần việc thật.
\textit{(If you want real ability, sooner or later you must go through the real work.)}
\end{baihoc}

\begin{ghinhoanh}
\textbf{Short English to remember:}

Shortcuts may save time now but often cost depth later.

\end{ghinhoanh}

\begin{tuhocnhanh}
1. Vì sao người ta hay tìm đường tắt? \textit{Why do people look for shortcuts?}

2. Đường tắt thường lấy đi điều gì? \textit{What do shortcuts often take away?}

3. Phần việc thật nào em đang muốn né? \textit{What real work are you trying to avoid?}
\end{tuhocnhanh}
\ngancach

\section{Luyện Tập Tạo Nên Kỹ Năng}
\begin{truyen}{Luyện Tập Tạo Nên Kỹ Năng}{Practice Builds Skill}
\chuhoa{K}{ỹ năng không đến chỉ vì em hiểu về nó. Nó đến khi em làm đi làm lại đủ nhiều.}
\textit{(Skill does not come because you understand it in theory. It comes when you repeat it enough.)}

Biết cách giải khác với giải được thật.
\textit{(Knowing the method is different from actually doing it.)}

Nói được quy tắc khác với dùng được quy tắc.
\textit{(Reciting a rule is different from using it well.)}

Luyện tập biến điều xa lạ thành quen tay, quen đầu.
\textit{(Practice turns the unfamiliar into something your hand and mind can handle.)}
\end{truyen}

\begin{baihoc}
Muốn giỏi hơn, em không chỉ cần biết thêm. Em còn cần làm thêm.
\textit{(To improve, you do not only need to know more. You also need to do more.)}
\end{baihoc}

\begin{ghinhoanh}
\textbf{Short English to remember:}

Skill grows through repeated doing.

\end{ghinhoanh}

\begin{tuhocnhanh}
1. Vì sao hiểu chưa đủ để thành kỹ năng? \textit{Why is understanding alone not enough for skill?}

2. Em đang biết mà chưa làm được ở chỗ nào? \textit{Where do you know but still cannot do?}

3. Em cần lặp lại điều gì nhiều hơn? \textit{What do you need to repeat more often?}
\end{tuhocnhanh}
\ngancach

\section{Học Là Hành Trình Dài}
\begin{truyen}{Học Là Hành Trình Dài}{Learning Is a Long Journey}
\chuhoa{N}{ếu em chỉ chịu cố khi thấy kết quả gần, em sẽ mệt với những đoạn đường dài của việc học.}
\textit{(If you only work hard when the result seems near, you will struggle with the long stretches of learning.)}

Nhưng học vốn là hành trình dài.
\textit{(But learning is naturally a long journey.)}

Có lúc rõ, có lúc mờ; có lúc tiến nhanh, có lúc đứng lại.
\textit{(Some phases are clear, others unclear; some move fast, others feel still.)}

Biết nó dài giúp em bớt hoảng khi một đoạn đường không đẹp như mong muốn.
\textit{(Knowing it is long helps you panic less when one segment does not look good.)}
\end{truyen}

\begin{baihoc}
Người đi đường dài tốt thường không phải người hưng phấn nhất, mà là người bền nhất.
\textit{(Those who travel long roads well are often not the most excited, but the most steady.)}
\end{baihoc}

\begin{ghinhoanh}
\textbf{Short English to remember:}

Long journeys reward steady hearts.

\end{ghinhoanh}

\begin{tuhocnhanh}
1. Vì sao biết việc học là hành trình dài lại quan trọng? \textit{Why is it important to see learning as a long journey?}

2. Đoạn nào của hành trình học làm em mệt nhất? \textit{What part of the learning journey exhausts you most?}

3. Em cần bền hơn ở đâu? \textit{Where do you need more steadiness?}
\end{tuhocnhanh}
\ngancach

\section{Bền Bỉ Là Sức Mạnh}
\begin{truyen}{Bền Bỉ Là Sức Mạnh}{Perseverance Is Strength}
\chuhoa{B}{ền bỉ không ồn, nhưng rất mạnh. Nó là khả năng ở lại với điều khó đủ lâu để điều khó dần bớt khó.}
\textit{(Perseverance is quiet but powerful. It is the ability to stay with difficulty long enough for difficulty to lessen.)}

Người bền không phải lúc nào cũng vui.
\textit{(Persevering people are not always cheerful.)}

Họ chỉ là không tự rút lui quá sớm.
\textit{(They simply do not withdraw too early.)}

Trong học tập, sức mạnh này quý hơn nhiều lời hứa lớn.
\textit{(In learning, this strength is worth more than many grand promises.)}
\end{truyen}

\begin{baihoc}
Bền bỉ là sức mạnh giúp con người đi qua những đoạn mà cảm hứng không còn cứu nổi.
\textit{(Perseverance is the strength that carries people through stretches where inspiration can no longer save them.)}
\end{baihoc}

\begin{ghinhoanh}
\textbf{Short English to remember:}

Perseverance keeps moving when inspiration fades.

\end{ghinhoanh}

\begin{tuhocnhanh}
1. Bền bỉ khác hưng phấn ở đâu? \textit{How is perseverance different from excitement?}

2. Vì sao bền bỉ là sức mạnh thật? \textit{Why is perseverance real strength?}

3. Em đang cần bền thêm với điều gì? \textit{What do you need to stay with longer?}
\end{tuhocnhanh}
\ngancach

\section{Đừng Bỏ Cuộc Giữa Đường}
\begin{truyen}{Đừng Bỏ Cuộc Giữa Đường}{Do Not Quit Midway}
\chuhoa{N}{hiều người bỏ ngay đoạn khó nhất, trong khi chỉ cần đi thêm một quãng nữa là mọi thứ có thể bắt đầu sáng hơn.}
\textit{(Many people quit at the hardest stretch, while just a little more distance might have made things clearer.)}

Giữa đường là lúc mệt nhất vì em đã tốn công mà chưa tới nơi.
\textit{(The middle is often the hardest because you have spent effort but not yet arrived.)}

Chính vì vậy, bỏ cuộc giữa đường rất hấp dẫn.
\textit{(That is why quitting midway is so tempting.)}

Nhưng rất nhiều điều đáng giá chỉ đến với người đi qua được cái đoạn lưng chừng ấy.
\textit{(But many worthwhile things come only to those who survive the in-between stretch.)}
\end{truyen}

\begin{baihoc}
Đừng để sự mệt của đoạn giữa làm em tưởng đích đến không còn giá trị.
\textit{(Do not let the exhaustion of the middle make you think the destination has no value.)}
\end{baihoc}

\begin{ghinhoanh}
\textbf{Short English to remember:}

The middle feels hardest because you are not there yet.

\end{ghinhoanh}

\begin{tuhocnhanh}
1. Vì sao đoạn giữa đường dễ làm người ta bỏ nhất? \textit{Why is the middle stretch where people most want to quit?}

2. Điều gì thường đến sau đoạn lưng chừng khó nhất? \textit{What often comes after the hardest middle stretch?}

3. Em đang đứng giữa chặng nào? \textit{What journey are you in the middle of right now?}
\end{tuhocnhanh}
\ngancach

\section{Thành Công Cần Thời Gian}
\begin{truyen}{Thành Công Cần Thời Gian}{Success Needs Time}
\chuhoa{N}{hiều học sinh muốn cố ít mà thấy thành quả nhiều. Nhưng đời ít khi vận hành như vậy.}
\textit{(Many students want little effort and quick reward. Life rarely works that way.)}

Thành công thường chậm hơn mong muốn của mình.
\textit{(Success is often slower than we want.)}

Nó đòi thời gian để tích lũy kỹ năng, va chạm, sửa sai và lớn lên.
\textit{(It requires time to accumulate skill, experience, correction, and growth.)}

Hiểu điều đó giúp em bớt nản khi hôm nay làm đúng mà ngày mai vẫn chưa thấy hoa nở.
\textit{(Understanding this helps you become less discouraged when good work today does not bloom tomorrow.)}
\end{truyen}

\begin{baihoc}
Không phải thứ đến chậm là vô ích. Nhiều điều quý chỉ thành hình khi được nuôi đủ lâu.
\textit{(What arrives slowly is not useless. Many valuable things take time to form.)}
\end{baihoc}

\begin{ghinhoanh}
\textbf{Short English to remember:}

Good growth often appears later than you want.

\end{ghinhoanh}

\begin{tuhocnhanh}
1. Vì sao học sinh hay sốt ruột với thành công? \textit{Why are students impatient with success?}

2. Thành công cần tích lũy những gì? \textit{What does success require to accumulate?}

3. Em có đang đòi kết quả quá sớm không? \textit{Are you demanding results too early?}
\end{tuhocnhanh}
\ngancach

\section{Người Cố Gắng Luôn Tiến Bộ}
\begin{truyen}{Người Cố Gắng Luôn Tiến Bộ}{Those Who Keep Trying Usually Improve}
\chuhoa{T}{iến bộ không phải lúc nào cũng lớn và rõ. Nhưng người cố gắng thật thường hiếm khi đứng yên mãi.}
\textit{(Improvement is not always large and visible. But those who truly keep trying rarely stay unchanged forever.)}

Có thể em chưa thấy ngay.
\textit{(You may not see it immediately.)}

Có thể người khác cũng chưa nhận ra.
\textit{(Others may not notice it either.)}

Nhưng nếu em đang học đều hơn, bình tĩnh hơn, bớt sợ hơn, hiểu chắc hơn, thì em đã tiến rồi.
\textit{(But if you study more steadily, feel calmer, fear less, and understand more solidly, then you have already moved forward.)}
\end{truyen}

\begin{baihoc}
Nỗ lực thật hiếm khi vô ích. Nó luôn để lại dấu vết trong con người em, dù sớm hay muộn.
\textit{(Real effort is rarely wasted. It always leaves traces in you, sooner or later.)}
\end{baihoc}

\begin{ghinhoanh}
\textbf{Short English to remember:}

Real effort leaves real change.

\end{ghinhoanh}

\begin{tuhocnhanh}
1. Vì sao tiến bộ không phải lúc nào cũng dễ thấy? \textit{Why is progress not always easy to see?}

2. Những dấu hiệu tiến bộ nào ngoài điểm số em có thể nhận ra? \textit{What signs of progress beyond scores can you notice?}

3. Em đã thay đổi điều gì nhờ cố gắng? \textit{What has changed in you through effort?}
\end{tuhocnhanh}
\ngancach
